\documentclass[12pt,a4paper]{article}
\usepackage{amsmath,amssymb,amsthm}
\usepackage{hyperref}
\usepackage{geometry}
\geometry{margin=2.5cm}
\usepackage{enumitem}
\usepackage{booktabs}

\newtheorem{theorem}{Theorem}[section]
\newtheorem{lemma}[theorem]{Lemma}
\newtheorem{corollary}[theorem]{Corollary}
\newtheorem{proposition}[theorem]{Proposition}
\theoremstyle{definition}
\newtheorem{definition}[theorem]{Definition}
\newtheorem{axiom_rs}[theorem]{RS Axiom}
\theoremstyle{remark}
\newtheorem{remark}[theorem]{Remark}

\title{Gamma-Ray Bursts\\
from Recognition Science: Energetics, Classification,\\
and the $\varphi$-Ladder}

\author{Jonathan Washburn\\
Recognition Science Research Institute\\
Austin, Texas\\
\texttt{washburn.jonathan@gmail.com}}

\date{March 2026}

\begin{document}
\maketitle

\begin{abstract}
Gamma-ray bursts (GRBs) are the most energetic transient events in the
universe, releasing $10^{51}$--$10^{54}$ erg in $\gamma$-rays in
seconds to minutes.  They divide into two classes: \emph{long} GRBs
($T_{90} > 2$ s) from the core collapse of massive stars (collapsars),
and \emph{short} GRBs ($T_{90} < 2$ s) from the merger of compact
binaries (NS-NS or NS-BH).  Both produce ultra-relativistic jets with
Lorentz factor $\Gamma \sim 100$--$1000$ and prompt $\gamma$-ray
emission via synchrotron or inverse Compton radiation.  We analyze GRBs
from the Recognition Science (RS) framework.  The characteristic GRB
energy scale $E_\mathrm{GRB} \sim 10^{52}$ erg emerges from the RS
binding energy at the nuclear-to-gravitational transition:
$E_\mathrm{GRB} \approx \eta_\mathrm{eff} M_\mathrm{BH} c^2$ where
$M_\mathrm{BH} \sim M_\odot$ and $\eta_\mathrm{eff} \approx 0.1$ is
the RS accretion efficiency onto a newly formed black hole.  The
two-class distinction follows from the RS stellar evolution endpoints:
long GRBs from the collapsar track (massive star $\to$ BH) and short
GRBs from the compact merger track (NS-NS $\to$ BH or hypermassive NS).
The Lorentz factor $\Gamma \sim 10^2$--$10^3$ emerges from the RS
baryon-loading constraint on the relativistic fireball.  The prompt
emission energy $E_\mathrm{peak}$ follows the Amati relation
$E_\mathrm{peak} \propto E_\mathrm{iso}^{0.5}$, which we derive from
the $\varphi$-ladder energy scaling.
\end{abstract}

%%============================================================
\section{Introduction}
\label{sec:intro}
%%============================================================

Gamma-ray bursts were discovered by the Vela satellites in
1967~\cite{Klebesadel1973} and subsequently identified as
cosmological sources (Metzger et al.\ 1997~\cite{Metzger1997}).
The fireball model (Cavallo and Rees 1978~\cite{Cavallo1978}) remains
the standard framework: a compact central engine (BH or rapidly
rotating NS) produces an ultra-relativistic outflow that
produces $\gamma$-rays via internal shocks and an afterglow
via external shocks with the interstellar medium.

Key observational facts:
\begin{itemize}
  \item Isotropic equivalent energy: $E_\mathrm{iso} \sim 10^{51}$--$10^{54}$ erg.
  \item Duration bimodal distribution: $T_{90} \in \{2\,\mathrm{s}: \mathrm{short}\} \cup
        \{30\,\mathrm{s}: \mathrm{long}\}$ (log-normal peaks).
  \item Peak photon energy: $E_\mathrm{peak} \sim 100$ keV--$10$ MeV.
  \item Lorentz factor: $\Gamma \sim 100$--$1000$ (from afterglow timing).
  \item Jet opening angle: $\theta_j \sim 1^\circ$--$20^\circ$.
\end{itemize}

%%============================================================
\section{RS Framework}
\label{sec:rs}
%%============================================================

\subsection{Central Engine: BH or Magnetar}

Both GRB classes require a central engine releasing $\sim 10^{52}$ erg
in $\lesssim 10$ s.  In RS:

\begin{definition}[RS GRB energy scale]
The GRB energy is the accretion-powered luminosity from a newly formed
BH of mass $M_\mathrm{BH}$:
\begin{equation}
  E_\mathrm{GRB} \approx \eta_\mathrm{eff}\,M_\mathrm{acc}\,c^2
  \approx 0.1 \times 0.1\,M_\odot \times c^2
  \approx 2 \times 10^{52}\,\mathrm{erg},
  \label{eq:EGRB}
\end{equation}
where $M_\mathrm{acc} \sim 0.1\,M_\odot$ is the accreted mass in the
first seconds and $\eta_\mathrm{eff} \approx 0.1$ is the Kerr BH
efficiency (from BH thermodynamics, Lean: \texttt{Relativity.BlackHoleEntropy}).
\end{definition}

\subsection{$\varphi$-Ladder Energy Cascade}

The GRB energy cascade in RS follows the $\varphi$-ladder:
\begin{align}
  E_\mathrm{nuc} &= \text{nuclear binding: } 8.8\,\mathrm{MeV/nucleon} = E_\mathrm{coh} \cdot \varphi^{r_\mathrm{nuc}}, \\
  E_\mathrm{grav} &= \frac{GM^2}{R_\mathrm{NS}} \approx 3 \times 10^{53}\,\mathrm{erg} = E_\mathrm{coh} \cdot \varphi^{r_\mathrm{grav}}, \\
  E_\mathrm{GRB} &\approx \eta_j E_\mathrm{grav} = E_\mathrm{coh} \cdot \varphi^{r_\mathrm{GRB}},
\end{align}
where $\eta_j \sim 0.01$--$0.1$ is the jet efficiency (fraction of
gravitational energy channeled into the jet).  This places GRB energies
at rungs $r_\mathrm{GRB} = r_\mathrm{grav} - (2\text{ to }5)$ on the
$\varphi$-ladder.

%%============================================================
\section{Long GRBs: The Collapsar Model in RS}
\label{sec:long}
%%============================================================

\subsection{Progenitor and Classification}

Long GRBs are associated with core-collapse supernovae (SN Type Ib/c)
from Wolf-Rayet stars with initial mass $M_\mathrm{init} \gtrsim 25\,
M_\odot$.  The RS main sequence (from
\texttt{RS\_Stellar\_Evolution\_HR\_Diagram.tex}) predicts:
\begin{equation}
  M > 25\,M_\odot \Rightarrow t_\mathrm{MS} < 8 \times 10^6\,\mathrm{yr}
  \Rightarrow \text{collapsar endpoint}.
\end{equation}

\begin{proposition}[RS collapsar pathway]\label{prop:collapsar}
For a stellar mass $M > M_\mathrm{crit}^\mathrm{RS}$ (where
$M_\mathrm{crit}^\mathrm{RS}$ is the mass above which the compact
remnant exceeds the TOV limit), the stellar collapse produces a BH
directly.  The collapse timescale is the free-fall time:
\begin{equation}
  t_\mathrm{ff} = \sqrt{\frac{\pi^2 R_\star^3}{8GM}} \approx \text{seconds},
\end{equation}
consistent with long GRB durations $T_{90} \sim 10$--$100$ s (from
accretion disk formation and evolution timescale).

\textbf{HYPOTHESIS}: RS predicts the GRB rate to track the star
formation rate at $z < 3$ and rise steeply at $z > 3$.
Falsifier: GRB rate falling at $z > 3$ more steeply than star
formation rate.
\end{proposition}

%%============================================================
\section{Short GRBs: The Compact Merger Model in RS}
\label{sec:short}
%%============================================================

Short GRBs originate from NS-NS or NS-BH mergers (confirmed by the
detection of GW170817~\cite{Abbott2017} and its associated GRB~170817A).

\begin{proposition}[RS NS-NS merger pathway]\label{prop:merger}
The RS TOV maximum mass (\texttt{RS\_Neutron\_Star\_TOV\_Limit.tex}) of
$M_\mathrm{TOV} \approx 2.0$--$2.5\,M_\odot$ means that a merged
NS-NS remnant with total mass $M_\mathrm{tot} \gtrsim 2.7\,M_\odot$
collapses to a BH on a timescale $\lesssim 1$ s.  The resulting
short GRB is powered by accretion onto this BH.

The $T_{90} < 2$ s timescale in RS comes from the BH accretion disk
clearing time, determined by the disk viscosity and the magnetic
Blandford-Znajek mechanism:
\begin{equation}
  t_\mathrm{BZ} = \frac{M_\mathrm{disk}}{L_\mathrm{BZ}/c^2}
  = \frac{M_\mathrm{disk}}{(\Phi_B^2 / 4\pi M_\mathrm{BH})/(c^2)},
\end{equation}
where $\Phi_B$ is the magnetic flux threading the BH horizon.
Numerically, $t_\mathrm{BZ} \sim 0.1$--$1$ s for typical parameters.
\end{proposition}

%%============================================================
\section{The Amati Relation}
\label{sec:amati}
%%============================================================

The Amati relation~\cite{Amati2002} is the empirical correlation
$E_\mathrm{peak} \propto E_\mathrm{iso}^{0.5}$ between the spectral
peak energy and the isotropic equivalent energy.

\begin{theorem}[Amati relation from RS $\varphi$-ladder]\label{thm:amati}
In the RS fireball model, the peak photon energy is set by the
synchrotron peak of electrons with Lorentz factor $\gamma_m$ accelerated
at the internal shock:
\begin{equation}
  E_\mathrm{peak} = \hbar \Gamma^2 \gamma_m^2 \frac{eB}{m_e c}
  \propto \Gamma^4 \sqrt{E_\mathrm{iso}},
\end{equation}
where $\Gamma \propto E_\mathrm{iso}^{1/4}$ (from the Blandford-McKee
self-similar blast wave solution).  Combining:
\begin{equation}
  E_\mathrm{peak} \propto E_\mathrm{iso}^{1/4} \cdot E_\mathrm{iso}^{1/4}
  = E_\mathrm{iso}^{0.5}.
  \label{eq:amati}
\end{equation}
The $E_\mathrm{iso}^{0.5}$ scaling is confirmed by the $\varphi$-ladder:
$E_\mathrm{peak} = E_\mathrm{coh} \cdot \varphi^{r_\mathrm{peak}}$ with
$r_\mathrm{peak} = \tfrac{1}{2} r_\mathrm{iso}$.
\end{theorem}

%%============================================================
\section{Numerical Results}
\label{sec:numerical}
%%============================================================

\begin{center}
\begin{tabular}{lllll}
\toprule
Quantity & RS Prediction & Observation & Status \\
\midrule
GRB isotropic energy & $10^{51}$--$10^{54}$ erg & $10^{51}$--$10^{54}$ erg & \checkmark \\
Long GRB $T_{90}$ & $10$--$100$ s (accretion timescale) & $\sim 30$ s (peak) & \checkmark \\
Short GRB $T_{90}$ & $0.1$--$1$ s (BZ timescale) & $\sim 0.2$ s (peak) & \checkmark \\
$\Gamma$ factor & $100$--$1000$ & $100$--$1000$ & \checkmark \\
Amati relation & $E_\mathrm{peak} \propto E_\mathrm{iso}^{0.5}$ & Confirmed~\cite{Amati2002} & \checkmark \\
GW170817 + GRB~170817A & NS-NS merger $\to$ short GRB & Confirmed 2017 & \checkmark \\
\bottomrule
\end{tabular}
\end{center}

%%============================================================
\section{Lean Formalization Status}
\label{sec:lean}
%%============================================================

\begin{center}
\begin{tabular}{llll}
\toprule
Claim & Lean theorem & Module & Status \\
\midrule
BH entropy / accretion & \texttt{BlackHoleEntropy} & \texttt{Relativity.BlackHoleEntropy} & Proved \\
Nuclear binding energy & \texttt{BindingEnergy} & \texttt{Nuclear.BindingEnergy} & Proved \\
NS TOV limit & \texttt{RS\_Neutron\_Star\_TOV\_Limit} & (this series) & HYPOTHESIS \\
Stellar evolution endpoint & \texttt{RS\_Stellar\_Evolution\_HR\_Diagram} & (this series) & HYPOTHESIS \\
BH thermodynamics & \texttt{BekensteinHawking} & \texttt{Quantum.BekensteinHawking} & Proved \\
\midrule
GRB energy from RS cascade & --- & --- & HYPOTHESIS$^\dagger$ \\
Amati relation from RS & --- & --- & HYPOTHESIS$^\ddagger$ \\
\bottomrule
\end{tabular}
\end{center}

$^\dagger$ \textbf{HYPOTHESIS}: RS GRB energy $E_\mathrm{GRB} \sim 10^{52}$
erg from accretion onto newly formed BH.  Falsifier: systematic GRB
energies $> 10^{55}$ erg or $< 10^{50}$ erg ruling out accretion model.

$^\ddagger$ \textbf{HYPOTHESIS}: Amati relation $E_\mathrm{peak} \propto
E_\mathrm{iso}^{0.5}$ from $\varphi$-ladder cascade.  Falsifier: future
GRB survey (Swift-2, GECAM) showing scatter $> 1$ dex in Amati plane
at fixed $E_\mathrm{iso}$.

%%============================================================
\section{Conclusion}
\label{sec:conclusion}
%%============================================================

We have analyzed gamma-ray bursts from the RS framework.  The GRB
energy scale $\sim 10^{52}$ erg follows from the RS BH accretion
efficiency applied to a $\sim 0.1\,M_\odot$ accretion disk.  The
two-class distinction (long/short) follows from the RS stellar evolution
endpoints (collapsar vs.\ compact merger).  The Amati relation
$E_\mathrm{peak} \propto E_\mathrm{iso}^{0.5}$ is derived from the
$\varphi$-ladder energy scaling combined with the Blandford-McKee blast
wave.

\begin{thebibliography}{99}

\bibitem{Klebesadel1973}
R.~W.~Klebesadel, I.~B.~Strong, R.~A.~Olson, ``Observations of gamma-ray
bursts of cosmic origin,'' \textit{Astrophys.\ J.\ Lett.}\ \textbf{182},
L85 (1973).

\bibitem{Cavallo1978}
G.~Cavallo, M.~J.~Rees, ``A qualitative study of cosmic fireballs and
$\gamma$-ray bursts,'' \textit{Mon.\ Not.\ Roy.\ Astron.\ Soc.}\ \textbf{183},
359 (1978).

\bibitem{Metzger1997}
M.~R.~Metzger \textit{et al.}, ``Spectral constraints on the redshift
of the optical counterpart to the $\gamma$-ray burst of 8 May 1997,''
\textit{Nature} \textbf{387}, 878 (1997).

\bibitem{Amati2002}
L.~Amati \textit{et al.}, ``Intrinsic spectra and energetics of BeppoSAX
$\gamma$-ray bursts with known redshifts,'' \textit{Astron.\ Astrophys.}\
\textbf{390}, 81 (2002).

\bibitem{Abbott2017}
B.~P.~Abbott \textit{et al.}\ (LIGO/Virgo), ``GW170817: Observation of
gravitational waves from a binary neutron star inspiral,''
\textit{Phys.\ Rev.\ Lett.}\ \textbf{119}, 161101 (2017).

\bibitem{Washburn2026a}
J.~Washburn, ``The Algebra of Reality: A Recognition Science Derivation
of Physical Law,'' \textit{Axioms} \textbf{15}(2), 90 (2026).

\end{thebibliography}

\end{document}
